\section{Training}
\label{sec:training}

During training with a small subset of data, we realized that we heavily differ from the original 100\% accuracy of the detector, as our dataset of 20 animals per breed was too small. Even though the Cat-API, from which we pulled our data, has a large dataset, the data is not properly labeled, and we could only obtain around 12 pictures per breed. 
As hand-labeling another eight pictures to reach the target of 20 pictures per breed takes too long, we avoided going forward with this and took our chance with further officially provided datasets. 

Here, we implemented the Oxford Dataset, which increased our pool by 
2,744 pictures of the animal breeds we needed. However, upon checking, 
we realized that \textit{Tabby} \& \textit{Tiger\_cat} are very similar, 
and \textit{Tabby} is not a breed but a fur pattern. In this case, 
we proceeded with \textit{Tiger\_cat}. Down the line, we recognized 
that we had too few pictures of \textit{Golden\_Retriever}, 
\textit{German\_Shepherd}, \textit{Siberian\_Husky}, \textit{Dalmatian}, 
and \textit{Rottweiler}, as these are not included in the Oxford Dataset. 
Due to this, we went further and implemented the Stanford Dogs Dataset, 
which has 150 pictures per specified breed and helped us to finalize our 
training dataset.
As we still struggled with tabby and tiger cat, we proceeded just for these breeds
to pull 30  Pictures from Unsplash.com and Pixabay.com. 
As we were still unsattisfied with the training results especially for these breeds,
we increased the density to 50 pictures per breed, which pushed the performance over the 80\% accuracy, which we held as our 
main goal.

During training we realized that our algorithm was still getting worse from time to time, due to this we verified our picture dataset and realized that dalmatian have been excluded for now 
of the dataset. As none of the api connections we used earlier had any dalamtian pictures, we donwloaded a file of 359 Dalmatian pictures and included this in our dataset. Through this after 10 Epochs of training we
had a big jump from 17 to 69\% accuracy.


learnings:
Batch size is to high -> lowering to 16
lowering pixel size to 224 instead of 384

Der Speicherverbrauch auf der GPU wächst nicht linear, sondern quadratisch mit der Bildgröße, da ein Bild zwei Dimensionen (Breite $\times$ Höhe) hat:Ein Bild mit der Standardgröße $224 \times 224$ hat 50.176 Pixel.Ein Bild mit deiner Größe $384 \times 384$ hat 147.456 Pixel.Das bedeutet: Nur weil du img_size von 224 auf 384 erhöht hast, muss deine Grafikkarte für jedes einzelne Bild fast die 3-fache Menge an Pixeln und Aktivierungen im VRAM abspeichern! Bei tieferen Netzwerken wie EfficientNet multipliziert sich dieser Speicherplatz durch die vielen Schichten des Modells extrem schnell auf.

Die batch_size = 32 bestimmt, wie viele dieser Bilder das neuronale Netz gleichzeitig (parallel) verarbeiten soll.Wenn PyTorch einen Vorwärtspass (Forward Pass) macht, muss es für alle 32 Bilder parallel die Aktivierungen jeder einzelnen Zwischenschicht im VRAM behalten, um danach im Rückwärtspass (Backward Pass) die Gradienten berechnen zu können.Bei batch_size = 32 und img_size = 384 versucht dein Skript grob gesagt, den Speicherplatz für $32 \times 147.456 \text{ Pixel} \times \text{Modell-Schichten}$ auf einmal zu reservieren. Das hat deine 7.6 GiB Grafikkarte gesprengt.Wenn du die batch_size auf 16 halbierst, muss die GPU in jedem Schritt nur noch halb so viele Zwischenergebnisse gleichzeitig speichern. Der Speicherverbrauch sinkt sofort um die Hälfte.

1. Das „Im-Kreis-Drehen“ (Sattelpunkte und Plateaus)
In der Realität dreht sich das Modell meistens nicht exakt im geometrischen Kreis, aber es fühlt sich so an: Wenn die Landschaft der Verlustfunktion (Loss-Landschaft) sehr flach ist (Plateau), sind die Gradienten ohnehin schon winzig. Multipliziert man das mit einer minimalen Lernrate, tendieren die Schritte gegen Null. Das Modell bewegt sich kaum noch vom Fleck und "verhungert" quasi auf dem Weg.

2. Die Falle der lokalen Minima


Eine sehr kleine Lernrate sorgt dafür, dass das Modell im allerersten kleinen "Tal" (lokales Minimum) gefangen bleibt, das ihm begegnet. Es hat schlichtweg nicht die "Energie" (Schrittweite), um über einen kleinen Hügel zu springen, um das eigentlich viel tiefere, bessere Tal (globales Minimum) dahinter zu finden.


Wenn du von bestimmten Rassen deutlich weniger Bilder hast als von anderen, tendiert das Modell sonst dazu, diese Minderheits-Klassen einfach zu ignorieren, um den Gesamt-Loss künstlich niedrig zu halten.Die eleganteste und ressourcenschonendste Methode in PyTorch ist es, der nn.CrossEntropyLoss()-Funktion einen Gewichtungs-Tensor (Weight-Tensor) zu übergeben.Die Mathematik dahinterWir berechnen die Gewichte für jede Klasse $i$ mit der Standardformel:$w_i = \frac{N_{total}}{C \cdot n_i}$Dabei ist $N_{total}$ die Gesamtzahl der Trainingsbilder, $C$ die Anzahl der Klassen (10) und $n_i$ die Anzahl der Bilder in der spezifischen Klasse.Häufige Rassen bekommen ein kleines Gewicht (z.B. $0.5$).Seltene Rassen bekommen ein hohes Gewicht (z.B. $3.0$). Ein Fehler bei einer seltenen Rasse wird das Modell also viel härter bestrafen.
Das ganze Prinzip basiert auf einem einfachen psychologischen Trick: Belohnung und Bestrafung. Wenn das Modell eine seltene Rasse falsch klassifiziert, wird es dafür viel härter bestraft als bei einer häufigen Rasse.Schritt 1: Das Zählen der BilderPythontrain_labels = []
for idx in train_idx:
    raw_label = int(base_dataset.data.iloc[idx, 1])
    label = raw_label if species == "cat" else raw_label - 10
    train_labels.append(label)
    
class_counts = Counter(train_labels)
Bevor das Training losgeht, läuft dieser Block einmal kurz durch deine Trainings-Indizes (train_idx). Er schaut in die Label-Spalte deines DataFrames und berechnet für Hunde das Label (indem er 10 abzieht, damit es wieder von 0 bis 9 läuft – genau wie in deinem Dataset).Der Counter erstellt daraus eine Liste, die dem Code sagt: "Hey, von Klasse 0 haben wir 500 Bilder, aber von Klasse 5 nur 20 Bilder!"Schritt 2: Die mathematische GewichtungHier berechnen wir für jede der 10 Klassen ihr individuelles Gewicht nach der Formel:$$w_i = \frac{N_{\text{total}}}{C \cdot n_i}$$$N_{\text{total}}$ = Gesamtzahl aller Trainingsbilder$C$ = Anzahl der Klassen (bei dir immer 10)$n_i$ = Anzahl der Bilder in dieser spezifischen Klasse $i$Pythonw = total_train_samples / (NUM_CLASSES * count)
Ein Rechenbeispiel, warum das funktioniert:Nehmen wir an, du hast insgesamt 1000 Bilder für 10 Katzenrassen.Rasse A (sehr häufig): Hat 500 Bilder.$$w_A = \frac{1000}{10 \cdot 500} = \frac{1000}{5000} = 0.2$$Rasse B (sehr selten): Hat nur 20 Bilder.$$w_B = \frac{1000}{10 \cdot 20} = \frac{1000}{200} = 5.0$$Die häufige Rasse bekommt also ein winziges Gewicht ($0.2$), die seltene Rasse ein riesiges Gewicht ($5.0$).Schritt 3: Das "Füttern" der Loss-FunktionPythonclass_weights_tensor = torch.tensor(weights, dtype=torch.float).to(device)
criterion = nn.CrossEntropyLoss(weight=class_weights_tensor)
PyTorch kann mit normalen Python-Listen während des Trainings nichts anfangen, weil das Ganze auf der GPU berechnet werden muss. Deshalb machen wir daraus einen torch.tensor und schieben ihn mit .to(device) auf die Grafikkarte.Diesen Tensor übergeben wir als weight an nn.CrossEntropyLoss().Schritt 4: Der Effekt während des TrainingsWenn dein Modell jetzt trainiert und ein Bild analysiert, läuft folgendes ab:Das Modell rät die Klasse.Es liegt falsch.Die Loss-Funktion berechnet den "Standard-Fehler" (sagen wir, der Fehlerwert ist mathematisch $1.0$).Jetzt kommen die Gewichte ins Spiel: * War das Bild von der häufigen Rasse A, rechnet der Loss: $1.0 \cdot 0.2 = \mathbf{0.2}$. Das Modell denkt sich: "Ah, nicht so wild."War das Bild von der seltenen Rasse B, rechnet der Loss: $1.0 \cdot 5.0 = \mathbf{5.0}$. Das ist ein gigantischer Fehler! Das Modell kriegt richtig "Ärger" (einen hohen Loss) und die Gewichte im neuronalen Netz verändern sich beim Backpropagation-Schritt drastisch, um diesen Fehler beim nächsten Mal unbedingt zu vermeiden.


Weiterer Versuch AI bilder zu generieren oder zu downloaden ist gescheitert bzgl. arbeitsaufwand... -> nicht vorhandener datensätze online


Ja, das ist absolut normal und gehört sogar zum Standard-Workflow im Machine Learning! Was du da beschreibst, nennt sich in der Praxis Checkpoint Saving (bzw. Early Stopping) in Kombination mit Hyperparameter-Tuning.

Dass Modelle mit mehr Epochen plötzlich wieder schlechter werden, ist ein altbekanntes Phänomen. Hier ist die Erklärung, warum das passiert und warum dein Vorgehen genau richtig ist:

Warum mehr Epochen nicht immer besser sind
Wenn du ein Modell länger trainierst, heißt das nicht automatisch, dass es schlauer wird. Oft passiert das Gegenteil:

1. Overfitting (Überanpassung)
Das ist der häufigste Grund. Ab einem gewissen Punkt lernt das Modell die Trainingsdaten einfach auswendig, anstatt die zugrundeliegenden Muster zu verstehen.

Die Folge: Die Accuracy auf den Trainingsdaten steigt weiter, aber die Accuracy auf den Validierungsdaten (die das Modell nicht kennt) stürzt ab.

2. Zu große Lernrate (Learning Rate)
Wenn deine Lernrate nicht mit der Zeit kleiner wird (Learning Rate Decay), kann es sein, dass das Modell am Ende des Trainings um das Optimum "herumspringt" oder sich komplett aus einem guten Zustand herauskatapultiert.

Wie Profis das lösen (Dein Ansatz automatisiert)
Genau das, was du händisch machst (den besten Run heraussuchen), automatisiert man im ML-Alltag über zwei Mechanismen:

Model Checkpointing
Du speicherst während des Trainings nach jeder Epoche das Modell ab – aber nur dann, wenn die Validierungs-Accuracy besser ist als der bisherige Bestwert.

Am Ende von 100 Epochen war vielleicht Epoche 42 die beste. Durch das Checkpointing wirfst du die verschlechterten Gewichte von Epoche 43–100 einfach weg und lädst den "Checkpoint" aus Epoche 42.

Early Stopping
Du sagst dem Code: "Wenn sich die Validierungs-Accuracy für 10 Epochen am Stück nicht verbessert, brich das Training sofort ab." Das spart Zeit und Rechenpower (und Slurm-Credits!).


wichtig: Bei jeder Epoche noch eingebaut ob neue train acc > als alte train acc und ob es nicht zu weit auseinanderdriftet um overfitting zu vermeiden. Wenn 

Overfitting vermeiden:
-> Viele Bilder
-> viele Epochen
-> Geringe learning curve
-> Dropouts
-> Cropmix